\begin{abstract}
Dynamic model-merging has emerged as a compelling paradigm for serving multiple task-specific expert adapters from a shared pre-trained backbone on resource-constrained devices, avoiding the prohibitive storage and memory footprint of independent model weights. However, existing dynamic model-merging schemes suffer from a severe and highly impractical trade-off. While static weight-merging suffers from ``heterogeneity collapse'' when processing mixed-task streams, state-of-the-art dynamic routing mechanisms bypass this by employing Micro-Batch Homogenization (MBH) to partition streams on-the-fly. Unfortunately, MBH requires up to $K$ sequential forward passes of the heavy base backbone, resulting in a linear latency penalty that is unacceptable for resource-constrained edge CPUs. Furthermore, traditional dynamic routers suffer from a temporal routing paradox: they require late-stage representation features to route, which forces them to execute the base model twice. 
To resolve these core deployment barriers and paradoxes, we propose \textbf{SPS-ZCA} (Single-Pass Sample-Wise Routing with Zero-Shot Centroid Alignment), a training-free, extremely fast dynamic model-merging framework. SPS-ZCA implements \textit{Single-Pass Activation-Space Dynamic Blending} (SPS) to blend expert adapter activations sample-wise inside a single, parallel forward pass of the shared base model, converting sequential $O(K)$ latency back to a constant $O(1)$ pass. Simultaneously, \textit{Zero-Shot Centroid Alignment} (ZCA) routes inputs using robust, head-independent task centroids pre-computed from a tiny calibration split in the pre-trained backbone's early representation space (Layer 3), bypassing noisy classification heads and resolving the routing paradox.
On mixed-task streaming benchmarks, SPS-ZCA recovers \textbf{100.0\%} of the expert ceiling (Joint Mean of \textbf{79.80\%}), outperforming prior non-parametric SOTA methods by \textbf{+3.66\%} absolute accuracy. Crucially, we honestly characterize the ``serving gap'': while raw out-of-the-box deep learning frameworks (like standard PyTorch) experience dynamic framework overheads that offset these benefits under uncompiled CPU execution, we show this is a software framework limitation. Using an analytical hardware-aware cost model, we project that an optimized compiler-level loop layout slashes projected analytical execution costs by \textbf{3.90$\times$} on highly heterogeneous batches (from 776.4 ms to 199.0 ms), maintaining a flat execution profile. More importantly, we demonstrate a \textbf{verified physical 1.17$\times$ wall-clock speedup} out of the box in uncompiled PyTorch at low batch scales ($B=16$), where sequential kernel dispatch overheads dominate. Finally, we introduce Intra-Task Dispersion Calibration (IDC) to resolve asymmetric manifold variance, and diagonal GMM coordinate density estimators to achieve a \textbf{95.2\%} true OOD task rejection rate, delivering a highly robust serving system.
\end{abstract}
